\chapter{Action de Polyakov, dualité T et compactification de l'état HMT dans les écrans de la théorie M}
\label{chap:teoria-m}

\section{Persistance étendue et surface d’univers}

Une corde physique est un objet unidimensionnel étendu dont l'évolution balaie
une surface d'univers. Chaque point de la corde possède une coordonnée interne
et chaque instant ajoute une coordonnée temporelle ; leur combinaison forme
le domaine sur lequel vivent la métrique induite, la tension et les champs de
fond. Une corde fermée possède une section spatiale circulaire. Une corde ouverte
possède des extrémités soumises à des conditions aux limites. Ses vibrations organisent
des états et sa tension convertit l'aire parcourue par la surface d'univers en
action.

Une brane de dimension \(p\) est un objet étendu dont l'histoire forme un
volume d'univers de dimension \(p+1\). Les branes reçoivent les extrémités des cordes,
transportent des champs sur leur volume et peuvent envelopper des cycles d'un espace
compact. Cette définition fixe trois données différentes : la dimension de l'objet,
la dimension de son histoire et les conditions qui gouvernent sa frontière. La
tension, les charges et le spectre appartiennent à la réalisation dynamique qui
complète ces données géométriques.

HMT atteint cette frontière à partir d'une architecture de routes persistantes. Une
route conserve l'orientation, le feuillet, la phase, le quotient, la retenue et la mémoire des
croisements qui l'ont prolongée. Le passage à un objet étendu assigne cette
généalogie à une section distribuée sur une coordonnée interne. La
composition des routes se réalise alors comme couture de segments ; la mémoire
de frontière enregistre les enroulements et les intersections ; l'action mesure la
persistance complète sur la surface parcourue. Cette application conserve l'
ordre causal : on construit d'abord l'état enrichi, puis on déclare
sa réalisation comme corde, brane ou champ.

L'action de Polyakov offre une forme précise pour le codomaine physique :
\[
 S_{\rm P}[X,h]
 =-\frac{T}{2}\int d^2\sigma\,
 \sqrt{-h}\,h^{ab}g_{MN}(X)
 \partial_aX^M\partial_bX^N.
\]
Ici, \(T\) est la tension, \(h_{ab}\) la métrique de la surface d'univers et
\(X^M\) l'application vers l'espace cible. Une réalisation HMT de corde doit
produire ces objets, leur espace des états et leurs symétries à partir d'une application
typée ; la structure de route fournit le domaine généalogique sur lequel
cette application agit.

\section{Écrans exacts en dimensions \(11\), \(10\), \(26\) et \(24\)}

Le premier écran part du code ternaire étendu
\(C_W=[12,6,6]_3\). Le poinçonnage d'une coordonnée produit le code
\(G_{11}=[11,6,5]_3\). Ses boules de Hamming de rayon deux ont pour cardinal
\[
 V_3(11,2)=1+22+220=243=3^5,
\]
et les \(3^6\) mots du code satisfont
\[
 3^6V_3(11,2)=3^{11}.
\]
Par conséquent, ces boules partitionnent exactement l'espace ambiant. L'entier
onze mesure ici la longueur d'un code parfait et conserve son type
combinatoire.

Un deuxième écran utilise le module réel de permutation \(V_{12}\) associé aux
douze sommets icosaédriques. Le supplémentaire du vecteur uniforme se
décompose en
\[
 V_{11}=V_3\oplus V_{3'}\oplus V_5.
\]
L'état dodécaphasique \(K\) polarise une droite \(\ell_K\subset V_3\) et
détermine le projecteur
\[
 P_{10}=P_{11}-\frac{u_Ku_K^{\mathsf T}}{\lVert u_K\rVert^2}.
\]
De celui-ci descendent les décompositions exactes
\[
 V_{11}=\ell_K\oplus V_{10},\qquad
 V_{10}=V_{3'}\oplus W_7(K),
\]
de dimensions \(11=1+3+7\) et \(10=3+7\). Le passage
\(12\to11\) élimine le mode uniforme ; le passage \(11\to10\) élimine la
direction marquée par \(K\). Une compactification physique ajoute une droite
temporelle, un réseau stable de rang six et le transport qui préserve les
trois directions étendues.

Le troisième écran part du réseau
\[
 L_{26}=\Lambda_{24}\oplus U\cong II_{25,1},
\]
où \(U\) est le plan hyperbolique entier. Pour un vecteur isotrope
primitif \(e_+\), la réduction
\[
 e_+^\perp/\mathbb Ze_+\cong\Lambda_{24}
\]
est exacte. La condition d'orthogonalité retire une direction et le quotient
retire la direction nulle contenue dans son orthogonal ; on obtient ainsi le rang
vingt-quatre. Code, module réel et réseau appartiennent à des catégories
distinctes. Leurs applications explicites permettent de relier leurs rangs en conservant la
provenance de chacun.

L'incidence entre les six faces et les huit sommets du cube apporte en outre
un écran de signature spatiale et temporelle. Si \(M\) est sa matrice d'
incidence, alors
\[
 MM^{\mathsf T}=12P_u+4P_-.
\]
Le quotient par le noyau possède un rang de quatre et se décompose en
\(1\oplus3\). En déclarant la droite uniforme comme direction temporelle, la forme
normalisée devient \(\operatorname{diag}(-1,1,1,1)\). L'incidence démontre
le rang et le rapport spectral ; l'orientation déclarée fixe la signature.

\section{Impulsion, enroulement et rayon réciproque}

Dans une coordonnée compacte de rayon \(R\), l'entier \(m\) enregistre l'impulsion et
l'entier \(w\) enregistre l'enroulement. La forme
\[
 E_R(m,w)=\frac{m^2}{R^2}+w^2R^2
\]
possède l'involution
\[
 \mathsf T(m,w,R)=(w,m,R^{-1}),
 \qquad
 E_R(m,w)=E_{R^{-1}}(w,m).
\]
Dans \(\ell^2(\mathbb Z^2)\), l'échange
\(U_T\lvert m,w\rangle=\lvert w,m\rangle\) fournit l'équivalence
unitaire
\[
 U_TH(R)U_T^{-1}=H(R^{-1}).
\]
C'est la résidence algébrique de la dualité de rayon : impulsion et
enroulement sont deux lectures réciproques d'un même secteur réticulaire. La
réalisation comme dualité T incorpore l'échelle de corde, les oscillateurs et
l'action sur l'espace physique des états. Au rayon autodual, les deux
lectures partagent la même échelle et l'échange agit au sein du même
spectre.

\section{Interface en dimension \(11\) : triplet spectral, supersymétrie et opérateur \(\mathcal Q^2=H\)}

La théorie M organise un codomaine physique à onze dimensions avec une métrique
\(g_{MN}\), une trois-forme \(C_{MNP}\), un gravitino \(\psi_M\), une action et des
transformations de supersymétrie. Le pont s'écrit comme une application
\[
 \mathfrak R_M:\mathcal X_{\rm HMT}^{\rm enr}
 \longrightarrow\mathcal X_{11}^{\rm phys}.
\]
Sa source conserve la phase, l'orientation, la couture et la mémoire ; son codomaine ajoute
les champs et la dynamique undécadimensionnelle. La compactification d'une
direction doit transporter ces champs vers un écran à dix dimensions,
préserver la limite de basse énergie et entrelacer les dualités déclarées.

Une compactification candidate est typée par un triplet
\((x_\infty,\kappa,\mathcal U)\). La section \(x_\infty\) représente une
persistance prolongeable ; \(\kappa\) la porte vers un espace de modules ;
\(\mathcal U\) transporte la métrique, l'orientation, le registre et la frontière. Son
admissibilité exige le prolongement à toute profondeur, la commutation avec la
dualité, la descente du spectre au quotient et la conservation des domaines et des
codomaines des observables. Ces conditions transforment une analogie de
rangs en un problème mathématique vérifiable.

Les cordes apparaissent après la compactification du secteur undécadimensionnel ; les
branes apparaissent comme des réalisations étendues couplées à la trois-forme et à
ses duales. La branche exceptionnelle qui atteint \(\Lambda_{24}\),
\(V^\natural\) et le Monstre conserve la même source enrichie par
un autre foncteur. L'interface dimensionnelle et l'interface exceptionnelle sont des prolongements
différents d'une généalogie commune, chacun ayant son propre codomaine.

\begin{resultado}
La perfection de \(G_{11}\), la polarisation \(11\to10\), la réduction
isotrope \(26\to24\), l'écran cubique de signature \((3,1)\) et l'
involution de rayon sont des résultats algébriques démontrés. Ensemble, ils forment la
résidence exacte des rangs, des quotients et des dualités sur laquelle se construit
l'interface vers les cordes, les branes et la théorie M.
\end{resultado}

\begin{lecturafisica}
Une histoire étendue conserve la position, l'enroulement, l'orientation et la
frontière dans une même persistance. La compactification détermine quelles
directions apparaissent comme extension macroscopique et lesquelles agissent comme degrés
internes. La corde lit cette persistance sur une surface d'univers ; la
brane la lit sur un volume d'univers ; la théorie M coordonne les deux lectures
au moyen de champs undécadimensionnels.
\end{lecturafisica}

\begin{frontera}
La réalisation est réfutée si l'application \(\mathfrak R_M\) altère les types des
écrans, si la dualité cesse d'agir unitairement sur le spectre, si
la compactification perd son prolongement ou si les champs ne satisfont pas l'action et la
supersymétrie. La clôture physique requiert une tension, des oscillateurs, des amplitudes,
des anomalies contrôlées et une limite gravitationnelle cohérente pour chaque
compactification déclarée.
\end{frontera}
